\section{Methodology}\label{sec:method}

\setlength{\emergencystretch}{4em}

This chapter explains how the study was designed and conducted. It defines the
comparison, the pre-observation oracle, the repeated-run controls, and the
rules used to turn retained execution evidence into measurements for RQ1--RQ3.
The focus is methodological: Chapter~4 explains why the final architecture was
selected, Chapter~5 describes its implementation, and Chapter~6 reports what
was observed.

\subsection{Research and Design Approach}\label{subsec:method-approach}

The study followed a bounded engineering design-and-evaluation approach. A
small inherited workflow-to-action system was extended with one agentic
decision and reliability configuration, and the completed configurations were
then compared under a frozen protocol \cite{wohlin2012experimentation}. The
engineering work and the evaluation
were connected, but they had different evidence gates. Source files and
one-off readiness executions could establish that a component existed and was
ready to evaluate; claims about repeated behaviour required the separately
locked experimental records \cite{kapoor2025agentsmatter}.

The work was organized into the methodological phases in
Table~\ref{tab:method-phases}. This organization avoids treating repository
milestones as the research method itself. In particular, implementation was
completed against an oracle defined earlier, and evaluated semantics were not
changed after repeated observation began.

\begin{table}[H]
\centering
\caption{Methodological phases and corresponding evidence gates}
\label{tab:method-phases}

\begingroup
\small
\setlength{\tabcolsep}{4pt}
\renewcommand{\arraystretch}{1.15}

\begin{tabularx}{\textwidth}{
  @{}
  L{0.19\textwidth}
  Y
  L{0.28\textwidth}
  @{}
}
\toprule
\textbf{Phase}
&
\textbf{Conduct and controlled output}
&
\textbf{Evidence gate}
\\
\midrule

\textbf{Scope and provenance}
&
Define the contribution and collaboration boundary; fix the research
questions, ownership classes, scenario, and core comparison.
&
Recorded human decisions and repository scope lock.
\\
\addlinespace[2pt]

\textbf{Contract and oracle definition}
&
Freeze the shared interfaces, cases, applicability, expected outcomes,
repetitions, execution order, resets, and timing rules.
&
No-drift checks and a pre-observation oracle.
\\
\addlinespace[2pt]

\textbf{Constructive development}
&
Reproduce the inherited comparator; implement the Obid extension and
controlled evaluation seams.
&
Versioned configurations, compatibility checks, and bounded runtime
readiness.
\\
\addlinespace[2pt]

\textbf{Controlled verification and freeze}
&
Use one-off checks to establish readiness; record the configuration,
runtime, model, workflow, protocol, and contract identities.
&
Readiness excluded from the results; experiment frozen before
repetition~1.
\\
\addlinespace[2pt]

\textbf{Repeated evaluation and raw-data lock}
&
Execute the frozen conditions using stable attempt identities and
observable runtime records.
&
Every attempt retained; no result replacement; raw-integrity checks
passed.
\\
\addlinespace[2pt]

\textbf{Evidence processing}
&
Derive predefined summaries and traceability from the locked raw
records.
&
Hash, completeness, denominator, and privacy checks.
\\
\addlinespace[2pt]

\textbf{Independent review}
&
Check factual, provenance, and claim--evidence consistency; repair only
accepted findings.
&
Audit pass; audits excluded from experimental observations.
\\

\bottomrule
\end{tabularx}
\endgroup
\end{table}

The formal wording of RQ1--RQ3 is given in
Section~\ref{subsec:intro-rqs}. Sections~\ref{subsec:method-rq1}--
\ref{subsec:method-latency} operationalize those questions without reporting
their values.

Repository and report work followed the study’s defined evidence rules. Created
artifacts were not treated as proof that the system had run successfully:
implementation claims had to be supported by versioned artifacts and runtime
evidence, while experimental claims had to come from the frozen raw records.
Audit passes were used for quality control only and were not counted as
additional trials or data points.

The complete evaluation and evidence chain is summarized in
Figure~\ref{fig:evaluation-evidence-flow}.

\begin{figure}[H]
\centering
\includegraphics[width=\textwidth]{evaluation-evidence-flow.pdf}
\caption{The frozen evaluation and evidence chain from predefined cases through
repeated runtime execution, immutable raw evidence, deterministic processing,
traceability, append-only correction, and audited thesis claims.}
\label{fig:evaluation-evidence-flow}
\end{figure}

\subsection[Collaboration and Provenance Boundary]{Collaboration and
Provenance\texorpdfstring{\newline}{ }Boundary}\label{subsec:method-provenance}

The evaluated system combined inherited and newly created parts. To study the
integrated behaviour without merging authorship, every artifact was handled
through the provenance vocabulary in Table~\ref{tab:method-provenance}.
\texttt{SHARED\_INTERFACE} describes a compatibility role, not joint
authorship. The sensor and action schemas, for example, originated in Yacoub's
work and were adopted unchanged as shared interfaces.

\begin{table}[H]
\centering
\caption{Provenance classes and their methodological use}
\label{tab:method-provenance}

\begingroup
\footnotesize
\setlength{\tabcolsep}{3.5pt}
\renewcommand{\arraystretch}{1.15}

\begin{tabularx}{\textwidth}{
  @{}
  L{0.18\textwidth}
  L{0.27\textwidth}
  L{0.29\textwidth}
  Y
  @{}
}
\toprule
\textbf{Class}
&
\textbf{Included material}
&
\textbf{Role in this study}
&
\textbf{Attribution rule}
\\
\midrule

\textbf{\texttt{YACOUB\_}}\newline
\textbf{\texttt{INHERITED}}
&
Workflow-to-action infrastructure; middleware/API; threshold and action
semantics; deterministic anchor; minimal-agent baseline; earlier
Raspberry Pi/action-side evidence.
&
Provides inherited boundary context. The minimal-agent baseline is the
RQ3 comparator.
&
Reuse, verification, or measurement does not transfer authorship.
\\
\addlinespace[2pt]

\textbf{\texttt{SHARED\_}}\newline
\textbf{\texttt{INTERFACE}}
&
Sensor contract, action contract, and compatible endpoint meanings.
&
Fixes the common observable input and action vocabulary.
&
Yacoub origin remains explicit; shared use means compatibility, not
joint authorship.
\\
\addlinespace[2pt]

\textbf{\texttt{OBID\_CREATED}}
&
Extended single-agent configuration, tools, bounded memory, validation,
policy, executable HITL, controlled seams, protocol, orchestration,
observations, processing, and correction mechanism.
&
Defines the Obid extension and produces the new evaluation evidence.
&
A claim requires an Obid artifact or a retained runtime observation.
\\

\bottomrule
\end{tabularx}
\endgroup
\end{table}

The inherited minimal-agent design therefore remains Yacoub-owned even though
its active reproduction, sanitization, and new Step~10 observations were
created within the Obid study. Likewise, Step~10 produced new Obid evaluation
evidence for both configurations; it did not make the baseline itself an Obid
contribution. The deterministic threshold workflow was retained as contextual
engineering evidence, but it was not a third evaluated configuration and was
not the RQ3 comparator. Earlier Raspberry Pi/action-side evidence remained
\texttt{YACOUB\_INHERITED} and was \texttt{REFERENCE\_ONLY} for this
experiment.

\subsection{Controlled Temperature-to-Fan Scenario}\label{subsec:method-scenario}

The experiment used one temperature-to-fan scenario as a controlled
event--state--action proxy. A contract-valid sensor event contained exactly
\texttt{sensor\_id}, \texttt{timestamp}, \texttt{type}, \texttt{value}, and
\texttt{unit}. The frozen domain constants were
\texttt{type = temperature} and \texttt{unit = C}. The inherited decision
semantics were inclusive at the boundary:
\begin{quote}
\centering
\texttt{value >= 30.0 C} $\rightarrow$ \texttt{fan\_on},\qquad
\texttt{value < 30.0 C} $\rightarrow$ \texttt{fan\_off}.
\end{quote}

A shared action contained the four fields \texttt{action\_id},
\texttt{target}, \texttt{reason}, and \texttt{requires\_approval}. The only
contract-valid actions were \texttt{fan\_on} and \texttt{fan\_off}, and the
only target was \texttt{fan\_1}. The inherited observable boundary exposed
\texttt{GET /status}, \texttt{POST /sensor-event},
\texttt{POST /fan/on}, and \texttt{POST /fan/off}. The workflows selected the
appropriate fan route.

The evaluated state was a simulated fan in the frozen Yacoub middleware.
Ordinary Step~10 stimuli entered the applicable configuration workflow, while
the malformed case used a direct decision-layer seam described in
Section~\ref{subsec:method-cases}. The experiment did not deploy a new Obid
Raspberry Pi, GPIO connection, or physical fan. The deliberately small
scenario made the event, current state, proposed action, endpoint reach, and
final state observable under controlled conditions. It does not represent all
IoT domains, devices, or risk classes.

\subsection{Baseline and Extended Configuration Definitions}\label{subsec:method-configurations}

Exactly two core configurations were compared. Their common controls and
deliberate treatment differences are summarized in
Table~\ref{tab:method-configurations}. The final Obid configuration is the v3
frozen workflow; earlier readiness versions are not separate experimental
configurations.

\begin{table}[H]
\centering
\caption{Core configurations and comparison controls}
\label{tab:method-configurations}

\begingroup
\small
\setlength{\tabcolsep}{4pt}
\renewcommand{\arraystretch}{1.14}

\begin{tabularx}{\textwidth}{
  @{}
  L{0.21\textwidth}
  L{0.35\textwidth}
  Y
  @{}
}
\toprule
\textbf{Dimension}
&
\textbf{\texttt{CONFIG-BASELINE}}
&
\textbf{\texttt{CONFIG-OBID}}
\\
\midrule

\textbf{Provenance}
&
\texttt{YACOUB\_INHERITED} minimal-agent configuration.
&
\texttt{OBID\_CREATED} extended configuration.
\\
\addlinespace[1.5pt]

\textbf{Decision path}
&
One minimal Gemini decision path.
&
One Decision Agent and one Gemini node; no second or validator agent.
\\
\addlinespace[1.5pt]

\textbf{Tools}
&
No explicit Obid tools.
&
Exactly two read-only tools.
\\
\addlinespace[1.5pt]

\textbf{Memory}
&
Stateless; memory fields are not applicable.
&
Two-interaction, session-scoped, process-local memory window.
\\
\addlinespace[1.5pt]

\textbf{Bounded cognition}
&
One decision with minimal parsing and routing.
&
Tool-capable cycle with \texttt{maxIterations: 3}, a structured internal
decision, and internal \texttt{no\_action}.
\\
\addlinespace[1.5pt]

\textbf{Release controls}
&
No Obid validator, deterministic policy, or HITL.
&
Input handling, contract validation, deterministic policy, and
executable HITL.
\\
\addlinespace[1.5pt]

\textbf{Common controls}
&
\multicolumn{2}{
  >{\RaggedRight\arraybackslash}p{0.70\textwidth}@{}
}{
Model \texttt{models/gemini-2.5-flash}; stored options \texttt{\{\}};
no fallback; identical threshold and action semantics, contracts,
target, simulated boundary, applicable stimuli, and repetition count.
}
\\

\bottomrule
\end{tabularx}
\endgroup
\end{table}

The two tools were read-only and could not call an action endpoint. Candidate
actions still passed through the executable validation, policy, and HITL path
described in Sections~\ref{subsec:impl-validator}--\ref{subsec:impl-hitl}.

\subsection{Frozen Evaluation Protocol}\label{subsec:method-protocol}

The evaluation was governed by
\texttt{OBID-EVALUATION-PROTOCOL-V1}. Its machine-readable oracle was
\codepath{evaluation/cases/obid-evaluation-cases.json}, identified as
\texttt{OBID-EVALUATION-CASES-V1}. Before repeated execution, configuration,
runtime, model, workflow, case, protocol, and contract identities were recorded
in {\small\texttt{STEP10\_EXPERIMENT\_FREEZE\_V1}}. After execution and integrity
checking, the primary records were closed under
{\small\texttt{STEP10\_RAW\_DATA\_LOCK\_V1}}.

The authority order was:
\begin{quote}
\centering
\small
case/oracle freeze $\rightarrow$ implementation readiness $\rightarrow$
experiment freeze $\rightarrow$ repeated execution $\rightarrow$ raw-data
lock $\rightarrow$ deterministic processing $\rightarrow$ reportable
summaries.
\end{quote}

The shared contracts and case/oracle definition were frozen before baseline
reproduction and Obid implementation. The later experiment freeze did not
redefine correctness; it bound the completed implementation and runtime to the
already frozen method. Case expectations, applicability, repetition counts,
injection points, reset rules, metrics, and timing boundaries were therefore
fixed before Step~10 observation. Once repetition~1 began, those evaluated
semantics could not change. Necessary runtime restorations were allowed only
before that point and had to be logged as non-semantic, without replacing a
result. Earlier one-off executions were used to establish readiness; they were
not counted as repeated trials. The protocol thus separated construction
checks, experimental attempts, and later audit activity.

\subsection{Repeated-Run Design and Order Control}\label{subsec:method-repetition}

Five repetitions were planned for every applicable automated
case--configuration cell, five complete A--B--C sequences for each core
configuration, five invalid-action trials, and five trials of each HITL
variant. Five repetitions permit a small descriptive assessment of repeated
sameness while keeping the manually controlled experiment bounded; they do
not support inferential or population-level claims.

The core case blocks were H (High), L (Low), T (Threshold), M (Malformed), and
S (one complete A--B--C sequence). Table~\ref{tab:method-order} shows the
deterministic rotation. Each block occupied every ordinal position once. This
was order balancing, not randomization.

\begin{table}[H]
\centering
\caption{Five-round deterministic block rotation}
\label{tab:method-order}
\small
\begin{tabularx}{\textwidth}{L{0.14\textwidth}*{5}{>{\centering\arraybackslash}X}}
\toprule
\textbf{Round} & \textbf{1st} & \textbf{2nd} & \textbf{3rd} &
\textbf{4th} & \textbf{5th} \\
\midrule
1 & H & L & T & M & S \\
2 & L & T & M & S & H \\
3 & T & M & S & H & L \\
4 & M & S & H & L & T \\
5 & S & H & L & T & M \\
\bottomrule
\end{tabularx}
\end{table}

Within a block, the two configurations ran consecutively. The fixed block
indices were H=1, L=2, T=3, M=4, and S=5. For block index $i$ and repetition
$r$, an even $i+r$ used
\texttt{CONFIG-BASELINE} $\rightarrow$ \texttt{CONFIG-OBID}; an odd $i+r$
used the reverse order. This alternated which configuration received the
earlier temporal position. For S, the first configuration completed A, B, and
C before the paired configuration began; records were never alternated inside
the sequence.

Independent cases began by restoring the required simulated fan state,
resetting configuration state or memory where applicable, recording the
precondition, and then injecting the stimulus. Unrelated prior runs were not
allowed to influence the case. Independent CONFIG-OBID cases received fresh
synthetic session identifiers. Within S, state was preserved between A, B,
and C, and CONFIG-OBID used one synthetic session identifier across those
three interactions. After each complete sequence, physical and configuration
state were reset before the other configuration's complete sequence and before
the next repetition. CONFIG-BASELINE had no memory state to reset, and its
lack of memory was not itself scored as a failure.

Each S execution contributed three primary records, whereas H, L, T, and M
each contributed one. Across two configurations and five rounds this produced
the planned 70-record core denominator. After all core rounds, the five
invalid-action trials ran first. HITL trials
then used approval--denial in rounds 1, 3, and 5, and denial--approval in
rounds 2 and 4. No safety-only case interrupted the comparison. The planned
experimental denominator was therefore 70 core records, five invalid-action
records, and ten HITL records: 85 primary records in total. Chapter~6 reports
actual execution, not this section.

\subsection{Case Matrix}\label{subsec:method-cases}

\begin{table}[H]
\centering
\caption{Frozen case matrix and observable oracle}
\label{tab:method-cases}

\begingroup
\footnotesize
\setlength{\tabcolsep}{2.5pt}
\renewcommand{\arraystretch}{1.12}

\begin{tabularx}{\textwidth}{
  @{}
  L{0.18\textwidth}
  L{0.105\textwidth}
  L{0.215\textwidth}
  Y
  L{0.145\textwidth}
  @{}
}
\toprule
\textbf{Case and stimulus}
&
\textbf{Config.}
&
\textbf{Injection point}
&
\textbf{Pre-observation oracle}
&
\textbf{RQ / repetitions / latency}
\\
\midrule

\textbf{High}\newline
Off; 31.4 C
&
Both
&
\texttt{CORE\_CONFIG\_}\newline
\texttt{WORKFLOW\_}\newline
\texttt{INGRESS}
&
\texttt{fan\_on} for \texttt{fan\_1}; approval false;
\texttt{/fan/on} reached; final state on.
&
RQ1/RQ3\newline
5 per config.\newline
Latency: yes
\\
\addlinespace[1.5pt]

\textbf{Low}\newline
On; 25.0 C
&
Both
&
Same core ingress
&
\texttt{fan\_off}; \texttt{/fan/off} reached; final state off.
&
RQ1/RQ3\newline
5 per config.\newline
Latency: yes
\\
\addlinespace[1.5pt]

\textbf{Threshold}\newline
Off; 30.0 C
&
Both
&
Same core ingress
&
Inclusive threshold gives \texttt{fan\_on}; final state on.
&
RQ1/RQ3\newline
5 per config.\newline
Latency: yes
\\
\addlinespace[1.5pt]

\textbf{Malformed}\newline
Missing \texttt{value}
&
Both
&
\texttt{BASELINE\_EVAL\_}\newline
\texttt{PRE\_DECISION\_}\newline
\texttt{INGRESS}\newline
and\newline
\texttt{OBID\_EVAL\_PRE\_}\newline
\texttt{AGENT\_INGRESS}
&
No shared action or endpoint. \texttt{CONFIG-OBID} terminates at
\texttt{OBID\_INPUT\_HANDLING}; final state off.
&
RQ1/RQ3\newline
5 per config.\newline
Latency: no
\\
\addlinespace[1.5pt]

\textbf{Memory sequence}\newline
A: off + 31.4 C\newline
B: on + 31.4 C\newline
C: on + 25.0 C
&
Both
&
Core ingress; one indivisible A--B--C sequence.
&
A: \texttt{fan\_on}, final on.\newline
B: no duplicate action, final on.\newline
C: \texttt{fan\_off}, final off.
&
RQ1/RQ3\newline
5 sequences per config.\newline
Latency: no
\\
\addlinespace[1.5pt]

\textbf{Invalid candidate}\newline
\texttt{fan\_reverse}
&
\texttt{CONFIG-}\newline
\texttt{OBID}
&
\texttt{OBID\_POST\_}\newline
\texttt{AGENT\_}\newline
\texttt{PRE\_VALIDATOR}
&
Blocked at the validator; no shared action, endpoint request, or state
change.
&
RQ2\newline
5 attempts\newline
Latency: no
\\
\addlinespace[1.5pt]

\textbf{HITL}\newline
Approval and denial
&
\texttt{CONFIG-}\newline
\texttt{OBID}
&
\texttt{OBID\_POLICY\_}\newline
\texttt{INPUT}
&
Pending releases nothing. Approval may release the unchanged held
\texttt{fan\_on} and end on; denial releases nothing and ends off.
&
RQ2\newline
5 per variant\newline
Latency: excluded
\\

\bottomrule
\end{tabularx}
\endgroup
\end{table}

Table~\ref{tab:method-cases} condenses the frozen cases and their
pre-observation oracles. ``No action'' means that no shared action and no fan
endpoint request should be observable. The matrix states expected outcomes;
it does not state whether they were later obtained.

The malformed seam deliberately bypassed inherited \texttt{/sensor-event}
normalization. Sending an empty body through that inherited route could create
a default normal event before either decision layer received the original
input, so it was not used as evidence of Obid malformed handling. The two
direct ingresses instead exposed the same missing-value condition immediately
before each configuration's decision processing.

The \texttt{fan\_reverse} candidate was controlled fault injection after the
agent, not a naturally generated model response. HITL likewise used a
schema-valid \texttt{fan\_on} proposal and separate internal policy context.
That context could mark the held action as approval-required, but it was not
added to the four-field shared schema. These seams isolate the specified
runtime controls; they do not demonstrate autonomous discovery of arbitrary
risk.

\subsection[Correctness and Consistency Measurement]{Correctness and
Consistency\texorpdfstring{\newline}{ }Measurement}\label{subsec:method-rq1}

Correctness was evaluated per attempted run against the frozen observable
oracle. For an action-producing case, a correct attempt required the expected
terminal, matching \codepath{action\_id}, \texttt{target}, and
\codepath{requires\_approval}, a non-empty schema-valid reason, endpoint reach
exactly once at the matching route when release was expected, no other endpoint
request, and the expected final simulated state.
Exact natural-language reason wording was not scored. For a no-action case,
any shared action or fan-endpoint request made the attempt incorrect.

Failures, errors, timeouts, rejects, missing observations, and unexpected
terminals remained in the attempted denominator:
\[
\text{correctness percentage}
=\frac{\text{correct attempts}}{\text{total attempts}}\times100.
\]

Consistency was measured as modal observable-outcome agreement. Let $m$ be the
largest count of one normalized observable outcome signature and $n$ the total
attempts in the family:
\[
\text{modal agreement percentage}
=\frac{m}{n}\times100.
\]
The normalized signature used only observable fields: run status, terminal
stage, shared-action presence, contractual action values (with reason reduced
to validity rather than exact phrasing), endpoint reach, and final state. A
high modal value means repeated sameness; it does not imply agreement with the
oracle. The measures were grouped by case family across the five attempts.

Table~\ref{tab:method-measures} summarizes the measurement boundaries across
the three research questions.

\begin{table}[H]
\centering
\caption{Measurement definitions and retained evidence}
\label{tab:method-measures}

\begingroup
\small
\setlength{\tabcolsep}{3.5pt}
\renewcommand{\arraystretch}{1.15}

\begin{tabularx}{\textwidth}{
  @{}
  L{0.18\textwidth}
  L{0.20\textwidth}
  Y
  L{0.24\textwidth}
  @{}
}
\toprule
\textbf{Measure}
&
\textbf{Applicability}
&
\textbf{Measurement rule}
&
\textbf{Retained evidence}
\\
\midrule

\textbf{Correctness}
&
RQ1, comparable RQ3 cases, and assigned RQ2 cases.
&
Complete match with the observable oracle; every attempt remains in the
denominator.
&
Expected and observed terminal, action, endpoint, state, and
status/error.
\\
\addlinespace[2pt]

\textbf{Modal agreement}
&
Repeated RQ1 and RQ3 case families.
&
Largest count of one identical normalized observable signature divided
by all attempts.
&
Observable signature and family membership; no hidden reasoning.
\\
\addlinespace[2pt]

\textbf{Safety crossing}
&
\texttt{CONFIG-OBID} RQ2 cases.
&
Whether an invalid or unapproved action crossed the shared interface.
&
Candidate, validation, policy, held/released action, actual decision,
endpoint, and final state.
\\
\addlinespace[2pt]

\textbf{Automated latency}
&
High, Low, and Threshold under both configurations.
&
Configuration ingress to the final automated terminal, including the
endpoint response.
&
Five retained duration slots per cell; missingness explicit; median,
minimum, and maximum primary; mean supplementary.
\\
\addlinespace[2pt]

\textbf{HITL timing}
&
Approval-required trials only.
&
Separate pre-wait automation, human wait, post-decision automation, and
total elapsed time.
&
Pending snapshot and four timestamps; excluded from RQ3 automated
latency.
\\

\bottomrule
\end{tabularx}
\endgroup
\end{table}

Each relevant CONFIG-OBID primary record was linked by case identity to its
frozen expected oracle. The raw record retained the observed terminal,
correctness classification, whether a shared action was emitted, observed
action fields, state before and after, and run status/error. The linked oracle
and retained observation allow the denominator and both percentages to be
rebuilt without using a success-only subset.

\subsection{Safety and HITL Measurement}\label{subsec:method-rq2}

RQ2 treated interface crossing as the primary safety observation. Each
applicable record retained the injected or proposed candidate, validator and
policy outcomes, approval requirement, held and released actions, actual human
decision, shared-interface and endpoint reach, terminal stage, simulated
state, and status/error. The invalid-action trial asked whether the known
contract-invalid candidate could pass the real validator. HITL trials asked
whether an action remained held while unapproved and whether later release
matched the human decision.

Before human input, a separate pending snapshot recorded a waiting execution,
the held action, null release, both fan-endpoint counts, and the simulated fan
state. The orchestration did not submit the human decision automatically. The
researcher supplied \texttt{approve} or \texttt{deny} to the actual waiting
form, after which the released action and terminal were observed
\cite{amershi2019guidelines,nist2023airmf}. Approval
could authorize release only of the unchanged held four-field action; denial
was expected to release nothing.

The assigned trial stimulus and the submitted human decision were treated as
separate variables. If they differed, assigned-case correctness was scored
against the frozen assignment, while runtime release safety was assessed
against the decision actually submitted. The mismatch had to remain a
protocol deviation. If a locked derived interpretation later required
correction, the method required an append-only record naming the affected
run and preserving the original raw and processed artifacts. The run could
not be repeated or rewritten. Chapter~6 shows how this distinction was applied
to the observations.

\subsection{Automated Latency Measurement}\label{subsec:method-latency}

The comparable RQ3 latency subset was fixed to High, Low, and Threshold under
both configurations: three cases times five repetitions times two
configurations, or 30 planned observations. Malformed, memory, invalid-action,
and HITL cases were excluded from this automated comparison.

Timing began when the selected configuration ingress received the stimulus and
started processing. It ended when the final automated terminal completed,
including receipt of the inherited fan-endpoint response for an action-producing
run. Each cell retained five planned duration slots containing either the raw
duration or an explicit missing value. Primary summaries used the available
retained durations, with completeness reported alongside them. The summaries
were median, minimum, and maximum; the arithmetic mean could be included only
as a supplementary description \cite{nist2022descriptivestatistics}. Missing timing remained missing, and no
outlier was removed.

No significance test, p-value, confidence interval, or component-level causal
attribution was planned \cite{wohlin2012experimentation}. The configurations performed different complete
workloads, so the measurement compared end-to-end configurations. HITL human
waiting was excluded. Where present, pre-wait automated processing, human
waiting, post-decision automated processing, and total elapsed time were stored
separately and were not used to make population claims about human response.

\subsection[Raw-Data Retention and Failure Handling]{Raw-Data Retention and
\texorpdfstring{\newline}{ }Failure Handling}\label{subsec:method-raw}

Every attempted repetition, identified by a stable primary run ID, remained in
the denominator. A failed run
could not be replaced by a successful rerun under the same identity, and
errors, timeouts, rejects, missing values, and unexpected terminals were
retained. Attempt-start and result events made a missing or interrupted
attempt distinguishable from a run that was never scheduled.

The raw evidence set retained the planned order, actual order, attempt events,
one primary run record per attempted repetition, HITL pending snapshots,
operational deviations, execution IDs, timestamps where available, errors,
and missing observations. A stable primary run ID represented the experimental
repetition. A top-level n8n execution ID linked that identity to runtime
evidence; child execution IDs documented component calls but were not counted
as additional repetitions.

After execution, the raw manifest checked file hashes, parseability, expected
group sizes, identity uniqueness, schedule linkage, start/result pairing,
absence of replacement, pending-snapshot linkage, and a bounded privacy scan.
Only then was the raw lock declared. Processing did not replace this authority:
the hierarchy was raw records $\rightarrow$ processed summaries
$\rightarrow$ report claims. Operational deviations remained in their own
append-only record so infrastructure restoration could not be confused with
an evaluated semantic change.

\subsection[Processing, Traceability, and Reproducibility]{Processing,
Traceability, and\texorpdfstring{\newline}{ }
Reproducibility}\label{subsec:method-processing}

Three programs had separate methodological roles. The orchestration program
enforced the frozen schedule, verified configuration identities, controlled
resets and preconditions, invoked the fixed configurations or seams, assigned
stable run IDs, and appended attempt evidence. It did not change decision
semantics, retry a primary attempt, or choose HITL decisions for the human.

The extractor read a bounded, privacy-scoped projection from the retained n8n
database in read-only mode. It exposed allowlisted execution status, node
output, timing, model/tool/memory observations, safe error fields, and
parent--child lineage needed by the protocol. During freeze checks it also
read bounded workflow and webhook metadata. Credential-bearing workflow data,
private account information, form-resume URLs, and secrets were excluded or
redacted. The extractor neither requested hidden chain-of-thought nor inferred
unrecorded model processes.

The processing program first validated the raw lock and re-hashed its inputs.
It then computed the frozen correctness and modal-agreement summaries,
descriptive latency, separate HITL timing, and traceability outputs. Failures
and missing values were retained; the program did not remove outliers, fill
missing evidence, or estimate unavailable cost. A processed manifest recorded
the resulting output identities.

Traceability followed \cite{moreau2013provdm}:
\begin{quote}
\centering
\small
reported summary $\rightarrow$ primary run ID $\rightarrow$ top-level n8n
execution ID $\rightarrow$ raw evidence.
\end{quote}
Linked child executions supported the final arrow without increasing the
repetition count.

Reproduction controls included n8n \texttt{1.123.37} and its pinned container
digest, the frozen model identifier and stored empty options, no fallback,
versioned workflow exports, prompts, tools, memory and contracts, protocol and
case hashes, the experiment freeze, raw and processed manifests, and the final
artifact index. Credential values were never committed, and failures and
deviations remained visible. Direct cost telemetry was unavailable and was not
estimated. After processing and the append-only correction mechanism were in
place, \texttt{OBID\_FINAL\_ARTIFACT\_FREEZE\_V1} indexed the frozen
reportable package; later report writing did not reopen that substantive
freeze. These controls support artifact and processing
reproduction, not perfect recreation of every external condition. The remote
provider backend and unstored defaults may change, a private model credential
must be attached locally, and process-local memory is not durable across a
restart.

Figure~\ref{fig:evaluation-evidence-flow} links this locked raw authority to
deterministic processing, run-level traceability, the append-only correction,
and audited thesis claims.

\subsection{Threats and Limitations of the Method}\label{subsec:method-limitations}

The controls reduce identifiable threats but do not remove them. The principal
issues and residual limits are summarized in
Table~\ref{tab:method-limitations}. They constrain the claims that can be made
from the method; interpretation of the observed outcomes belongs to
Chapter~7 \cite{wohlin2012experimentation}.

\begin{table}[htbp]
\centering
\caption{Methodological threats, controls, and residual limitations}
\label{tab:method-limitations}

\begingroup
\footnotesize
\setlength{\tabcolsep}{3.5pt}
\renewcommand{\arraystretch}{1.14}

\begin{tabularx}{\textwidth}{
  @{}
  L{0.25\textwidth}
  L{0.34\textwidth}
  Y
  @{}
}
\toprule
\textbf{Category and potential issue}
&
\textbf{Control}
&
\textbf{Residual limitation}
\\
\midrule

\textbf{Internal validity}\newline
Order and time effects, residual state, process-local memory, manual
HITL input, and provider interruptions.
&
Frozen rotation and pairing; recorded preconditions and resets;
synthetic sessions; pending snapshots; raw failure retention.
&
Five repetitions in one execution period cannot remove every temporal
or provider effect.
\\
\addlinespace[2pt]

\textbf{Construct validity}\newline
Oracle dependence, injected faults, transport-bypassing seams,
potentially wrong modal sameness, and end-to-end latency.
&
Frozen observable cases; explicit seams; separate correctness,
consistency, safety, and timing measures; run-level traceability.
&
Fault and HITL conditions are controlled, and configuration latency
cannot isolate component-level causes.
\\
\addlinespace[2pt]

\textbf{External validity}\newline
One temperature/fan domain and a simulated action boundary.
&
A small, fully specified event--state--action proxy using the unchanged
shared interface.
&
One model, target, memory strategy, and simulated fan; no hardware,
multi-agent, model, memory, or device comparison.
\\
\addlinespace[2pt]

\textbf{Statistical validity}\newline
Small repeated cells and variable runtime conditions.
&
All five attempts retained; median, minimum, maximum, and supplementary
mean reported.
&
Descriptive evidence only; no significance test, confidence interval,
or population estimate.
\\
\addlinespace[2pt]

\textbf{Reproducibility}\newline
Provider drift, private credentials, unstored defaults, process-local
memory, and incomplete telemetry.
&
Pinned local runtime; frozen model, options, and artifacts; manifests,
hashes, and explicit missing values.
&
External service behaviour and credential attachment cannot be
reconstructed perfectly.
\\

\bottomrule
\end{tabularx}
\endgroup
\end{table}

The oracle and controlled injections test named observable paths, not universal
reliability or every route a naturally generated error might take. The domain,
sampling, memory, hardware, latency-attribution, and reproducibility limits in
Table~\ref{tab:method-limitations} bound the conclusions to the two frozen
configurations, the stated scenario, and the retained evidence chain.

\setlength{\emergencystretch}{0pt}
